\section{第七天}
\begin{enumerate}
  \item 坚持“三依三治”。
        \begin{enumerate}
          \item 坚持“三依三治”的概念和意义。
                \begin{enumerate}
                  \item 全面依法治国的系统性。
                        \begin{enumerate}
                          \item 全面依法治国是一个系统的工程。
                          \item 全面依法治国是国家治理领域一场广阔而深刻的变革，
                          \item 涉及改革发展稳定，内政国防外交，治党、治国治军等各个方面的问题，
                          \item 必须要立足全局和长远进行系统谋划。
                          \item 在全面依法治国过程当中，要坚持系统观念准确把握全面依法治国的整体工作安排，
                          \item 坚持依法治国、依法执政、依法行政共同推进，
                          \item 法治 国家、法治政府、法治社会一体建设。
                        \end{enumerate}
                  \item 三依。
                        \begin{enumerate}
                          \item 依法治国、依法执政、依法行政是一个有机整体，
                          \item 关键在于党要坚持依法执政、各级政府要坚持依法行政。
                        \end{enumerate}
                  \item 三治。
                        \begin{enumerate}
                          \item 法治国家、法治政府、法治社会三者各有侧重，相辅相成。
                          \item 法治国家是法治建设的目标，法治政府是建设法治国家的主体，法治社会是构筑法治国家的基础。
                        \end{enumerate}
                  \item 总结。全面推进依 法治国，必须着眼全局、统筹兼顾，在共同推进上着力，在一体建设上用劲。
                \end{enumerate}
          \item 法治国家、法治政府、法治社会（法治社会必须精确背下来，可考性极高极高）。
                \begin{enumerate}
                  \item 法治国家（目标）。
                        \begin{enumerate}
                          \item 法治国家是全面推进依法治国的根本目标。
                          \item 法治国家是指依法赋予、运行和制约国家权力、通过公正司法和严格执法来维护法律权威并实现人民权利的国家存在形式。
                        \end{enumerate}
                  \item 法治政府（主体）。
                        法治政府是政府依据宪法法律设立、政府权力法定、政府决策和行为严格依据法律程序进行并对其后果承担相应责任的政府。
                  \item 法治社会（基础）。
                        法治社会是社会依法治理、社会成员人人崇尚法治和信仰法治、社会组织依法自治、社会秩序在法治下和谐稳定的社会。
                \end{enumerate}
        \end{enumerate}
\end{enumerate}